% Part: sets-functions-relations
% Chapter: functions
% Section: functions-relations

\documentclass[../../../include/open-logic-section]{subfiles}

\begin{document}

\olfileid{sfr}{fun}{rel}

\olsection{સંબંધો તરીકે વિધેયો}

\begin{explain}
$A$ના !!{element}sને $B$ના !!{element}s પર મોકલતું
વિધેય સ્પષ્ટપણે $A$ અને~$B$ વચ્ચે સંબંધ વ્યાખ્યાયિત
કરે છે: $x$ અને $y$ વચ્ચે આ સંબંધ હોય જો અને માત્ર જો
$f(x) = y$ હોય. હકીકતમાં, દાખલા તરીકે ગણિતના પાયાના
ઘટકોની સંખ્યા ઘટાડવામાં રસ હોય તો, આપણે વિધેય~$f$ને
આ સંબંધ સાથે, એટલે કે જોડોના ગણ સાથે, \emph{એકરૂપ}
પણ ગણી શકીએ. ત્યારે પ્રશ્ન ઊભો થાય છે: કયા સંબંધો
આ રીતે વિધેયો વ્યાખ્યાયિત કરે છે?
\end{explain}

\begin{defn}[વિધેયનો આલેખ] ધારો કે $f\colon A \to B$
વિધેય છે. $f$નો \emph{આલેખ} એ નીચે પ્રમાણે વ્યાખ્યાયિત
સંબંધ $R_f \subseteq A \times B$ છે:
\[
R_f = \Setabs{\tuple{x,y}}{f(x) = y}.
\]
\end{defn}

\begin{explain}
ઘટકો દ્વારા નિર્ધારિત સમાનતાના સિદ્ધાંતથી વિધેયનો
આલેખ અનન્ય રીતે નક્કી થાય છે. વધુમાં, ગણો માટેનો
આ સિદ્ધાંત વિધેયો માટેના ગર્ભિત સમાનતાના સિદ્ધાંતને
તરત સમર્થન આપે છે: જો $f$ અને~$g$નો પ્રદેશ અને
સહપ્રદેશ એક જ હોય, તો બધી દલીલો માટેની કિંમતો
સરખી હોય ત્યારે તેઓ એક જ વિધેય હોય છે.

એ જ રીતે, સંબંધ “વિધેયાત્મક” હોય, તો તે કોઈ
વિધેયનો આલેખ છે.
\end{explain}

\begin{prop}\ollabel{prop:graph-function}
ધારો કે $R \subseteq A \times B$ એવો છે કે:
\begin{enumerate}
\item જો $Rxy$ અને $Rxz$ હોય, તો $y = z$; અને
\item દરેક $x \in A$ માટે કોઈક $y \in B$ એવો છે
કે $\tuple{x,
y} \in R$.
\end{enumerate}
તો $R$ એ $f(x) = y$ જો અને માત્ર જો $Rxy$
દ્વારા વ્યાખ્યાયિત વિધેય $f\colon A \to B$નો આલેખ છે.
\end{prop}

\begin{proof}
ધારો કે $Rxy$ થાય એવો $y$ છે. જો $Rxz$ થાય એવો
બીજો $z \neq
y$ હોત, તો $R$ પરની શરતનો ભંગ થાત.
આથી, $Rxy$ થાય એવો $y$ હોય, તો આ $y$ અનન્ય છે,
અને તેથી $f$ સુવ્યાખ્યાયિત છે. સ્પષ્ટપણે $R_f = R$.
\end{proof}

\begin{explain}
દરેક વિધેય $f\colon A \to B$નો આલેખ હોય છે, એટલે કે
$f(x) = y$ દ્વારા વ્યાખ્યાયિત $A
\times B$માં સમાયેલો સંબંધ. બીજી તરફ,
\olref{prop:graph-function}માં આપેલા ગુણધર્મો ધરાવતો
દરેક સંબંધ~$R
\subseteq A \times B$ કોઈ વિધેય~$f \colon A \to
B$નો આલેખ છે. વિધેયો અને તેમના આલેખો વચ્ચેના આ ગાઢ
સંબંધને કારણે આપણે વિધેયને તેના આલેખ તરીકે જ વિચારી શકીએ.
બીજા શબ્દોમાં, વિધેયોને અમુક સંબંધો સાથે, એટલે કે અમુક
બહુજોડોના ગણો સાથે, એકરૂપ ગણી શકાય.
\oliflabeldef{sfr:rel:ref:sec}{જોકે, નોંધો કે આ “એકરૂપતા”નો
આશય \olref[sfr][rel][ref]{sec}માં જણાવ્યો હતો તેવો છે:
તે વિધેયોના તત્વમીમાંસાકીય સ્વરૂપ વિશેનો દાવો નથી;
વિધેયોને અમુક ગણો તરીકે \emph{લેવાં} અનુકૂળ છે એવી
નોંધ છે. આ અનુકૂળતાનું એક કારણ એ છે કે આપણે}{આપણે}
હવે સંબંધો પર કરેલી ક્રિયાઓ જેવી ક્રિયાઓ વિધેયો પર
કરવાનું વિચારી શકીએ છીએ (જુઓ \olref[sfr][rel][ops]{sec}).
ખાસ કરીને:
\end{explain}
\sourcecorrection{OLFUN-004}{મૂળની \texttt{functions-relations.tex}, પંક્તિઓ
60--68માં આલેખને ``\(A\times B\) પરનો સંબંધ'' કહેવાયો છે. પ્રોજેક્ટની
પોતાની વ્યાખ્યા પ્રમાણે અહીં \(A\) અને \(B\) વચ્ચેનો, એટલે કે
\(A\times B\)માં સમાયેલો સંબંધ અભિપ્રેત છે; અનુવાદમાં એ પ્રકાર સ્પષ્ટ કર્યો છે.}

\begin{defn}\ollabel{defn:funimage}
ધારો કે $f \colon A \to B$ વિધેય છે અને $C\subseteq A$ છે.

$f$નું $C$ પરનું \emph{મર્યાદન} એ બધા $x \in C$ માટે
$(\funrestrictionto{f}{C})(x) = f(x)$ દ્વારા વ્યાખ્યાયિત
વિધેય~$\funrestrictionto{f}{C}\colon C \to B$ છે.
બીજા શબ્દોમાં, $\funrestrictionto{f}{C} = \Setabs{\tuple{x, y} \in R_f}{x \in
C}$.

$C$ પર $f$નું \emph{પ્રયોજન} એ $\funimage{f}{C} =
\Setabs{f(x)}{x \in C}$ છે. તેને $f$ હેઠળનું
$C$નું \emph{પ્રતિબિંબ} પણ કહીએ છીએ.
\end{defn}

\begin{explain}
આ વ્યાખ્યાઓ પરથી કોઈ પણ વિધેય~$f$ માટે
$\ran{f} =
\funimage{f}{\dom{f}}$ મળે છે.
\oliflabeldef{sfr:rel:ops:sec}{\olref[sfr][rel][ops]{sec}ની
વ્યાખ્યાઓ અને વિધેયોને સંબંધો સાથે એકરૂપ ગણવાની આપણી
રીત જોતાં, આ ખ્યાલો અનુરૂપ છે. અહીં વિધેયનું મર્યાદન
માત્ર પ્રદેશને C સુધી સીમિત કરે છે; અગાઉનું સંબંધ-મર્યાદન
બંને અવયવોને C સુધી સીમિત કરે છે. પરંતુ બે બીજી
ક્રિયાઓ---વ્યસ્ત અને સાપેક્ષ ગુણાકાર---માટે થોડી વધુ
વિગત જરૂરી છે. તે \olref[inv]{sec} અને
\olref[cmp]{sec}માં આપીશું.}{}
\end{explain}
\sourcecorrection{OLFUN-005}{મૂળની \texttt{functions-relations.tex}, પંક્તિઓ
78--100માં વિધેય-મર્યાદનની સ્પષ્ટ વ્યાખ્યા સાચી છે, પરંતુ તેને અગાઉના
બંને-અવયવવાળા સંબંધ-મર્યાદન સાથે ``ચોક્કસ'' સરખું કહેવાયું છે. અહીં
આ સામ્યની મર્યાદા સ્પષ્ટ કરી છે.}

\end{document}
